\section{Hyperparameter tuning}
\label{app:tuning}

All four methods share one PPO configuration, tuned for IPPO alone on
\texttt{simple\_spread} with a separate tuning seed (1000) and checked on
\texttt{simple\_adversary}. Selection used self-play return on the reserved
evaluation episodes. We varied the learning rate ($3{\times}10^{-4}$ to
$2{\times}10^{-3}$), epochs per update (5--15), minibatches (4--16), episodes per
update (32--256) and the entropy coefficient (0 or 0.01); the full table is in
\texttt{docs/tuning.md}. Each configuration was run with a single seed, so
differences of 1--2 return points are within seed-to-seed noise. The chosen
configuration (learning rate $10^{-3}$, 10 epochs, 8 minibatches, 128 episodes
per update, entropy 0.01) reached a training return still improving slowly at
$5\times10^{6}$ steps.

\section{Environment re-implementation}
\label{app:env}

Our batched environments follow the \texttt{mpe2} reference implementation of
\texttt{simple\_spread\_v3} (local reward ratio 0.5, agent radius 0.15, soft
contact forces) and \texttt{simple\_adversary\_v3} ($N=3$ good agents, one
adversary, shaped distance rewards), with discrete actions and 25-step
episodes. Resets reproduce the reference random-number draws exactly. The test
suite steps both implementations side by side under random actions from
several seeds and checks that observations agree to $10^{-5}$ and rewards to
$10^{-6}$ at every step. In \texttt{simple\_spread} the agents are silent, so
the two communication entries in each teammate block are always zero.

\section{Order sensitivity of the trained policies}
\label{app:sensitivity}

\begin{figure}[h]
\centering
\includegraphics[width=\linewidth]{figures/learning_curves.pdf}
\caption{Training return (top) and permutation sensitivity on the training
batch (bottom) over training; mean $\pm$ one standard error over 10 seeds. Deep
Sets is exactly invariant and is omitted from the bottom row. In
\texttt{simple\_adversary} the training return is measured against each run's
own co-trained adversary and is therefore not comparable across methods.}
\label{fig:curves}
\end{figure}

\section{Training and inference cost}
\label{app:cost}

\begin{table}[h]
\centering
\caption{Cost of one PPO update (128 episodes, 10 epochs of 8 minibatches) and
of one batched forward pass at inference time, measured in a single process on
one CPU thread on identical synthetic batches (median of 5 and 200 repeats).
\picr{} and IPPO deploy the same network. Order augmentation permutes inputs
while acting, so its update cost equals IPPO's by construction.}
\label{tab:cost}
\resizebox{\linewidth}{!}{\input{tables/overhead.tex}}
\end{table}
